% Manually-laid-out ontology overview.
% Module membership is encoded by node fill colour:
%   * Method module        → blue   (algofill)
%   * Training Data module → green  (trainfill)
%   * Benchmark module     → orange (benchfill)
%
% Layout invariants:
%   * Six rows defined by the central column (Implementation, MLIPMethod,
%     MLIPRun, TrainedModel, MaterialSystem, TrainingDataset).  All
%     central-column nodes share x = 5, the same minimum width, and a
%     uniform vertical spacing of 1.3 cm.
%   * Left column (Library, SimulationType, AtomicConfiguration) shares
%     west x = -1.1 (anchor=west).
%   * Right column (BenchmarkStudy, BenchmarkResult, AccuracyMetric,
%     MetricProperty) shares east x = 16 (anchor=east).
%   * Four leaf classes (DatasetProvenance, CoveredProperty, MetricType,
%     MetricProperty) sit on a new row at y = -0.8.
%   * Cross-module edges are dashed; the trainedWith and trainedUsing
%     shortcuts are dotted.
%   * Node and edge labels use the same DL macros as the main text
%     (\dlConcept and \dlRole), rendered in math sans-serif.
\begin{tikzpicture}[
    x=1cm, y=1cm,
    every node/.style={font=\footnotesize},
    methodnode/.style={draw, rounded corners=2pt, fill=algofill,
        minimum height=0.55cm, inner sep=3pt, align=center,
        font=\footnotesize, text height=1.6ex, text depth=0.4ex},
    trainnode/.style={draw, rounded corners=2pt, fill=trainfill,
        minimum height=0.55cm, inner sep=3pt, align=center,
        font=\footnotesize, text height=1.6ex, text depth=0.4ex},
    benchnode/.style={draw, rounded corners=2pt, fill=benchfill,
        minimum height=0.55cm, inner sep=3pt, align=center,
        font=\footnotesize, text height=1.6ex, text depth=0.4ex},
    centerwidth/.style={minimum width=2.5cm},
    edge/.style={->, >=stealth, thin},
    crossedge/.style={->, >=stealth, thin, dashed},
    shortcut/.style={->, >=stealth, thin, dotted},
    edgelabel/.style={font=\scriptsize, fill=white, inner sep=0.5pt,
        text height=1.4ex, text depth=0.3ex},
]
  \node[methodnode, centerwidth] (Implementation)  at (5, 7.0) {\ImplementationC};
  \node[methodnode, centerwidth] (MLIPMethod)      at (5, 5.7) {\MLIPMethodC};
  \node[methodnode, centerwidth] (MLIPRun)         at (5, 4.4) {\MLIPRunC};
  \node[methodnode, centerwidth] (TrainedModel)    at (5, 3.1) {\TrainedModelC};
  \node[trainnode,  centerwidth] (MaterialSystem)  at (5, 1.8) {\MaterialSystemC};
  \node[trainnode,  centerwidth] (TrainingDataset) at (5, 0.5) {\TrainingDatasetC};
  \node[methodnode, anchor=west] (Library)              at (-1.1, 7.0) {\LibraryC};
  \node[methodnode, anchor=west] (SimulationType)       at (-1.1, 5.7) {\SimulationTypeC};
  \node[trainnode,  anchor=west] (AtomicConfiguration)  at (-1.1, 1.8) {\AtomicConfigurationC};
  \node[methodnode, minimum width=3.2cm] (HyperparameterSetting) at (11.25, 3.1) {\HyperparameterSettingC};
  \node[benchnode,  anchor=east, minimum width=2.5cm] (BenchmarkStudy)  at (16.5, 3.1) {\BenchmarkStudyC};
  \node[benchnode,  anchor=east, minimum width=2.5cm] (BenchmarkResult) at (16.5, 1.8) {\BenchmarkResultC};
  \node[benchnode,  anchor=east, minimum width=2.5cm] (AccuracyMetric)  at (16.5, 0.5) {\AccuracyMetricC};
  \node[methodnode, minimum width=3.2cm] (FunctionalForm) at (11.25, 7.0) {\FunctionalFormC};
  \node[methodnode, minimum width=3.2cm] (LossFunction)   at (11.25, 5.7) {\LossFunctionC};
  \node[methodnode, minimum width=3.2cm] (Hyperparameter) at (11.25, 4.4) {\HyperparameterC};
  % Leaf classes — single row at y=-0.8
  \node[trainnode]  (DatasetProvenance) at (3,    -0.8) {\DatasetProvenanceC};
  \node[trainnode]  (CoveredProperty)   at (7,    -0.8) {\CoveredPropertyC};
  \node[benchnode]  (MetricType)        at (12.5, -0.8) {\MetricTypeC};
  \node[benchnode, anchor=east]  (MetricProperty)    at (16.5, -0.8) {\MetricPropertyC};

  % DFTSettings: same row as MetricType, centred horizontally in the gap
  % between CoveredProperty's east edge and MetricType's west edge.
  % DFTCalculation is one row above DFTSettings, sharing its x-centre.
  \node[trainnode] (DFTSettings)
    at ($(CoveredProperty.east)!0.5!(MetricType.west)$) {\DFTSettingsC};
  \node[trainnode] (DFTCalculation)
    at ($(CoveredProperty.east)!0.5!(MetricType.west) + (0, 1.3)$) {\DFTCalculationC};

  \draw[edge] (MLIPMethod) -- node[edgelabel] {\hasFunctionalFormR} (FunctionalForm.west);
  \draw[edge] (MLIPMethod) -- node[edgelabel] {\hasLossFunctionR} (LossFunction.west);
  \draw[edge] (MLIPMethod) -- node[edgelabel] {\hasHyperparameterR} (Hyperparameter.west);
  \draw[edge] (MLIPMethod) -- node[edgelabel] {\hasImplementationR} (Implementation);
  \draw[edge] (MLIPMethod) -- node[edgelabel] {\supportsSimulationR} (SimulationType);
  \draw[edge] (Implementation) -- node[edgelabel] {\implementedInR} (Library);
  \draw[edge] (HyperparameterSetting) -- node[edgelabel] {\forHyperparameterR} (Hyperparameter);
  \draw[edge] (MLIPRun) -- node[edgelabel] {\appliesMethodR} (MLIPMethod);
  \draw[edge] (MLIPRun) -- node[edgelabel] {\producesR} (TrainedModel.north);
  \draw[edge] (TrainingDataset) -- node[edgelabel] {\hasDFTCalculationR} (DFTCalculation);
  \draw[edge] (TrainingDataset) -- node[edgelabel] {\coversMaterialR} (MaterialSystem);
  \draw[edge] ($(TrainingDataset.north west)!2/3!(TrainingDataset.south west)$)
    to[out=180, in=270] node[edgelabel, pos=0.6] {\hasConfigurationR} (AtomicConfiguration.south);
  \draw[edge] (DFTCalculation) -- node[edgelabel] {\hasDFTSettingsR} (DFTSettings);
  \draw[edge] (BenchmarkStudy) -- node[edgelabel] {\hasResultR} (BenchmarkResult);
  \draw[edge] (BenchmarkResult) -- node[edgelabel] {\hasAccuracyMetricR} (AccuracyMetric);

  % Leaf-class edges
  \draw[edge] (TrainingDataset) -- node[edgelabel] {\datasetProvenanceR} (DatasetProvenance);
  \draw[edge] (TrainingDataset) -- node[edgelabel] {\coversPropertyR} (CoveredProperty);
  \draw[edge] (AccuracyMetric) -- node[edgelabel] {\metricTypeR} (MetricType);
  \draw[edge] (AccuracyMetric) -- node[edgelabel] {\metricPropertyR} (MetricProperty);

  % runsOn polyline: vertical middle segment runs through the exact
  % midpoint x of the corridor between AtomicConfiguration's east edge
  % and MaterialSystem's west edge.
  \coordinate (runsOnEnd) at ($(TrainingDataset.north west)!1/3!(TrainingDataset.south west)$);
  \coordinate (corridor)  at ($(AtomicConfiguration.east)!0.5!(MaterialSystem.west)$);
  \path let \p1=(corridor), \p2=(MLIPRun.west), \p3=(runsOnEnd) in
    coordinate (runsOnTopCorner) at (\x1, \y2)
    coordinate (runsOnBotCorner) at (\x1, \y3);
  \draw[crossedge, rounded corners=4pt]
    (MLIPRun.west) -- (runsOnTopCorner) --
    node[edgelabel] {\runsOnR}
    (runsOnBotCorner) -- (runsOnEnd);

  % evaluatesModel: leaves BenchmarkResult at 1/4 down its west side and
  % arrives at TrainedModel at 3/4 down its east side; both tangents
  % horizontal-westward (out=180, in=0).
  \draw[crossedge]
    ($(BenchmarkResult.north west)!1/4!(BenchmarkResult.south west)$)
    to[out=180, in=0]
    node[edgelabel, pos=0.5] {\evaluatesModelR}
    ($(TrainedModel.north east)!3/4!(TrainedModel.south east)$);
  \draw[crossedge] (BenchmarkResult) -- node[edgelabel] {\targetMaterialR} (MaterialSystem);

  % HyperparameterSetting incoming: solid from MLIPRun (underlying), dotted
  % shortcut from TrainedModel.  hasHyperparameterSetting leaves MLIPRun
  % heading east (out=0) and arrives at 1/4 down HyperparameterSetting's
  % west side, also heading east (in=180).
  \draw[edge] (MLIPRun.east) to[out=0, in=180]
    node[edgelabel, pos=0.5] {\hasHyperparameterSettingR}
    ($(HyperparameterSetting.north west)!1/4!(HyperparameterSetting.south west)$);
  \draw[shortcut] (TrainedModel.east) to
    node[edgelabel, pos=0.55] {\trainedUsingR}
    (HyperparameterSetting.west);

  % trainedWith shortcut: dotted arc on left of central column, swung
  % wider westward (looseness 2.5) so it does not overlap the runsOn
  % polyline.  Ends at 3/4 down MLIPMethod's west side.
  \draw[shortcut, bend left=70, looseness=2.5]
    (TrainedModel.west) to node[edgelabel, pos=0.5] {\trainedWithR}
    ($(MLIPMethod.north west)!3/4!(MLIPMethod.south west)$);
\end{tikzpicture}
